% Part: normal-modal-logic
% Chapter: axioms-systems
% Section: provability-properties

\documentclass[../../../include/open-logic-section]{subfiles}

\begin{document}

\olfileid{nml}{prf}{prp}

\olsection{Sipat-Sipat \usetoken{S}{derivability}}

\begin{prop}\ollabel{prop:derivabilityfacts}
  Upamane $\Sigma$ sistem modal lan $\Gamma$ himpunan !!{formula}s
  modal. Sipat-sipat ing ngisor iki lumaku:
  \begin{enumerate}
  \item\ollabel{prop:derivabilityfacts-monotonicity} \emph{Monotonisitas}: Yen
    $\Gamma \Proves[\Sigma] !A$ lan $\Gamma \subseteq \Delta$, mula
    $\Delta \Proves[\Sigma] !A$;
  \item\ollabel{prop:derivabilityfacts-reflexivity}
    \emph{Refleksivitas}: Yen $!A \in \Gamma$, mula $\Gamma
    \Proves[\Sigma] !A$;
  \item\ollabel{prop:derivabilityfacts-cut} \emph{Cut}: Yen $\Gamma
    \Proves[\Sigma] !A$ lan $\Delta \cup \{!A\} \Proves[\Sigma] !B$,
    mula $\Gamma \cup \Delta \Proves[\Sigma] !B$;
  \item\ollabel{prop:derivabilityfacts-deduction} \emph{Teorema deduksi}: $\Gamma \cup \{!B\}
    \Proves[\Sigma] !A$ yen lan mung yen $\Gamma \Proves[\Sigma] !B \lif
      !A$;
  \item \ollabel{prop:derivabilityfacts-ruleT} \emph{Aturan T}: Yen
    $\Gamma \Proves[\Sigma] !A_1$ lan \dots lan $\Gamma
    \Proves[\Sigma] !A_n$ lan $!A_1 \to (!A_2 \lif \cdots (!A_n \lif
    !B)\cdots)$ iku instansi tautologis, mula $\Gamma
    \Proves[\Sigma] !B$.
  \end{enumerate}
\end{prop}

Buktine dadi latihan kang gampang. Bagean
\olref{prop:derivabilityfacts-ruleT} saka \olref{prop:derivabilityfacts}
nuduhake, upamane, yen $\Gamma \Proves[\Sigma] !A \lor !B$
lan $\Gamma \Proves[\Sigma] \lnot !A$, mula $\Gamma \Proves[\Sigma]
!B$. Sabanjure, kita nulis $\Gamma, !A \Proves[\Sigma] !B$
minangka cekakan saka $\Gamma \cup \{ !A\} \Proves[\Sigma] !B$.

\begin{defn}
  Himpunan $\Gamma$ diarani \emph{katutup sacara deduktif} relatif marang
  sistem $\Sigma$ yen lan mung yen $\Gamma \Proves[\Sigma] !A$
  ngimplikasikake $!A \in \Gamma$.
\end{defn}

\end{document}
